\documentclass[12pt]{article}

\usepackage[letterpaper, margin=1in]{geometry}
\usepackage[spanish]{babel}
\usepackage[utf8]{inputenc}
\usepackage[T1]{fontenc}
\usepackage{setspace}
\usepackage{graphicx}
\usepackage{hyperref}
\usepackage{csquotes}
\usepackage{titlesec}

\usepackage[
    backend=biber,
    style=apa
]{biblatex}

\addbibresource{references.bib}

\onehalfspacing

\title{
    \textbf{Efecto de la familiaridad conductual en agentes conversacionales sobre la percepción del usuario}\\
    \large Propuesta de investigacion
}

\author{
    Luis Carlos Quesada Rodríguez \\
    Universidad de Costa Rica \\
    Escuela de Ciencias de la Computación e Informática
}

\date{\today}

\begin{document}

\maketitle

\begin{abstract}
Resumen breve de la propuesta. Debe explicar el problema, el agente propuesto,
el contexto de uso y el aporte esperado de la investigación.
\end{abstract}


\section{Definición del Problema}

Los agentes virtuales son generales

\subsection{Relevancia del problema}

Explique por qué el problema es importante desde una perspectiva social,
educativa, técnica o científica.

\subsection{Población afectada}

Indique a quién afecta el problema. Por ejemplo: estudiantes, docentes,
personas adultas mayores, pacientes, trabajadores, usuarios de una plataforma, etc.

\subsection{Limitaciones de las soluciones actuales}

Describa qué soluciones existen actualmente y cuáles son sus limitaciones.
Puede mencionar aspectos como accesibilidad, personalización, interacción,
costo, escalabilidad, falta de empatía, poca adaptabilidad, entre otros.







\section{Revisión de Literatura}

\subsection{Familiaridad como Señal Cognitiva}

La familiaridad ha sido ampliamente estudiada en psicología cognitiva como un componente fundamental del reconocimiento. En modelos de memoria de doble proceso, se distingue entre familiaridad y recuerdo, donde la familiaridad corresponde a un proceso rápido y automático que genera una sensación de conocimiento sin recuperar información contextual específica. A diferencia del recuerdo, que implica recuperación explícita de detalles, la familiaridad opera como una señal subjetiva de reconocimiento.

Investigaciones posteriores han ampliado esta perspectiva, mostrando que la familiaridad no depende exclusivamente de la memoria previa, sino que puede emerger a partir de la facilidad de procesamiento de un estímulo, fenómeno conocido como \textit{processing fluency}. Estímulos que son más fáciles de procesar tienden a ser percibidos como más familiares, incluso en ausencia de exposición previa. Esto sugiere que la familiaridad no es una medida directa de conocimiento, sino una inferencia basada en señales cognitivas.

Como consecuencia, la familiaridad puede influir en la forma en que los individuos evalúan información, favoreciendo juicios heurísticos y reduciendo el procesamiento analítico. Este carácter inferencial convierte a la familiaridad en un mecanismo relevante para entender cómo los usuarios interpretan y evalúan agentes virtuales.

\subsection{Familiaridad en la Percepción de Agentes Virtuales}

En el contexto de interacción humano-agente, la familiaridad ha sido asociada con efectos en la percepción del usuario, particularmente en dimensiones como la confianza, la naturalidad y la atribución de conocimiento. Estudios en agentes conversacionales han mostrado que los agentes percibidos como familiares tienden a ser considerados más confiables y más conocedores, incluso cuando su capacidad real no difiere de la de otros agentes.

Este fenómeno sugiere la existencia de una dimensión epistémica de la familiaridad, en la cual los usuarios atribuyen conocimiento al agente basándose en la sensación de familiaridad, más que en evidencia objetiva. Desde una perspectiva cognitiva, esto puede explicarse por la tendencia a depender de representaciones previas cuando se percibe familiaridad, reduciendo la necesidad de evaluar críticamente la información presentada.

No obstante, gran parte de la literatura ha abordado la familiaridad en agentes a través de la replicación de identidad, por ejemplo mediante \textit{digital doubles} o avatares de personas conocidas. Estos enfoques combinan múltiples factores, incluyendo apariencia, voz e identidad, lo que dificulta aislar los mecanismos específicos que generan los efectos observados.

\subsection{Antropomorfismo como Estrategia de Generación de Familiaridad}

Una estrategia común para inducir respuestas sociales en agentes virtuales ha sido el uso de antropomorfismo, entendido como la atribución de características humanas a entidades no humanas. Este enfoque se basa en la tendencia humana a interpretar agentes utilizando modelos sociales previamente aprendidos, lo cual puede aumentar la percepción de presencia social, credibilidad y engagement.

Sin embargo, la efectividad del antropomorfismo depende de la coherencia entre las señales presentadas por el agente. Investigaciones han mostrado que el incremento en el nivel de humanización puede generar expectativas más altas por parte del usuario, las cuales, si no son satisfechas, resultan en evaluaciones negativas. Este fenómeno se relaciona con el efecto conocido como \textit{uncanny valley}, donde representaciones altamente realistas pero imperfectas generan incomodidad y rechazo.

Estos resultados sugieren que los enfoques basados en humanización son sensibles a inconsistencias entre apariencia y comportamiento, lo que limita su robustez como mecanismo para inducir familiaridad.

\subsection{Alineación Conductual como Mecanismo Alternativo}

Más allá del antropomorfismo, investigaciones en psicología social y lingüística han mostrado que los humanos tienden a alinear su comportamiento durante la interacción, tanto en dimensiones no verbales como verbales. Este fenómeno incluye la imitación de gestos, así como la alineación en el uso del lenguaje, conocida como \textit{linguistic style matching}.

La alineación lingüística se manifiesta en la adaptación de elementos como vocabulario, estructura sintáctica y estilo comunicativo, y ha sido asociada con el aumento de rapport, confianza y conexión social. De manera similar, teorías como la \textit{Communication Accommodation Theory} proponen que los individuos ajustan su forma de comunicarse para reducir distancia social y mejorar la interacción.

A diferencia de los enfoques basados en identidad, estos mecanismos sugieren que la familiaridad puede emerger a partir de patrones de comportamiento compartidos, sin necesidad de replicar características perceptuales de una persona específica.

\subsection{Hacia la Familiaridad Conductual}

Los hallazgos anteriores permiten conceptualizar la familiaridad como un fenómeno que puede ser inducido mediante la alineación conductual. En este contexto, se propone entender la familiaridad conductual como la percepción de similitud entre los patrones de interacción de un agente y las expectativas del usuario sobre una persona conocida.

Dado que la familiaridad puede surgir a partir de la fluidez del procesamiento, es plausible que la alineación conductual facilite la interacción, aumentando la facilidad de procesamiento y, en consecuencia, generando una sensación de familiaridad. Sin embargo, la literatura actual no ha explorado de manera directa si la imitación conductual parcial, en ausencia de señales explícitas de identidad, es suficiente para inducir este efecto.

\subsection{Evaluación de la Percepción en Agentes Virtuales}

La evaluación de agentes virtuales se ha abordado principalmente mediante metodologías centradas en el usuario, en las cuales los participantes interactúan con agentes bajo diferentes condiciones y posteriormente reportan sus percepciones a través de cuestionarios estructurados. Este enfoque es consistente con prácticas establecidas en interacción humano-computadora, donde constructos como confianza, naturalidad o percepción de conocimiento son inherentemente subjetivos.

Estudios previos han utilizado escalas tipo Likert para medir estos constructos, permitiendo comparar condiciones experimentales de manera controlada. Este tipo de medición es particularmente adecuado en el contexto de la familiaridad, dado que esta se manifiesta como una experiencia subjetiva y no puede ser observada directamente mediante métricas objetivas del sistema.

\subsection{Brecha de Investigación}

En conjunto, la literatura muestra que la familiaridad influye en la percepción del usuario y que puede ser inducida mediante distintos mecanismos, principalmente a través de antropomorfismo y replicación de identidad. Sin embargo, estos enfoques presentan limitaciones importantes, incluyendo la dificultad de aislar variables y la sensibilidad a inconsistencias perceptuales.

Por otro lado, aunque existe evidencia de que la alineación conductual influye en la interacción, no se ha explorado si este mecanismo puede inducir familiaridad en ausencia de señales explícitas de identidad, ni cómo esto afecta la percepción de conocimiento previo por parte del usuario.

Este trabajo aborda esta brecha al investigar si la imitación conductual de una persona conocida puede generar percepciones de familiaridad y conocimiento en interacciones conversacionales breves, utilizando un enfoque centrado en el comportamiento en lugar de la apariencia o identidad.








\section{Preguntas de Investigación}

Las preguntas de investigación deben estar centradas en el agente como objeto de estudio.
\subsection{RQ1}

\textbf{RQ1:} ¿En qué medida la incorporación de comportamientos de imitación de una persona conocida en un Agente inteligente virtual influye en la percepción de naturalidad, cercanía y conocimiento previo por parte del usuario durante una conversación breve?

\subsection{RQ2}

\textbf{RQ2:} ¿En qué medida un Agente inteligente virtual con comportamientos de imitación de una persona conocida influye en las respuestas emocionales reportadas por los usuarios, en comparación con un agente sin dichos comportamientos?

\subsection{RQ3}

\textbf{RQ3:} ¿En qué medida un Agente inteligente virtual con comportamientos de imitación de una persona conocida es percibido como poseedor de conocimiento previo sobre el usuario, en comparación con un agente sin dichos comportamientos?

\subsection{RQ4 Opcional}

\textbf{RQ4:} ¿En qué medida un Agente inteligente virtual con comportamientos de imitación es percibido como similar a la persona conocida que intenta representar por parte del usuario?

\section{Objetivos del Agente}

\subsection{Rol del agente}

Describa qué papel cumple el agente.

Por ejemplo:

\begin{itemize}
    \item Asistente.
    \item Tutor.
    \item Guía.
    \item Compañero virtual.
    \item Entrenador.
    \item Facilitador de información.
\end{itemize}

Explique también cómo se comporta: si responde preguntas, da recomendaciones,
monitorea progreso, ofrece retroalimentación, guía actividades, etc.

\subsection{Contexto de uso}

Explique en qué situación interactúa el usuario con el agente.

Incluya:

\begin{itemize}
    \item Usuario objetivo.
    \item Lugar o entorno de uso.
    \item Tipo de interacción.
    \item Frecuencia esperada de uso.
    \item Necesidades principales del usuario.
\end{itemize}

\section{Arquitectura Técnica Preliminar}

Esta sección debe describir las posibles herramientas, tecnologías y decisiones técnicas
del agente propuesto.

\subsection{Descripción general de la arquitectura}

Incluya aquí una descripción general del sistema.

Ejemplo:

\begin{quote}
El sistema estará compuesto por una interfaz de usuario, un módulo de procesamiento
del lenguaje natural, un modelo de inteligencia artificial, una base de conocimiento,
servicios de voz opcionales y una representación visual del agente.
\end{quote}

\subsection{Diagrama preliminar}

Puede insertar una imagen del diagrama con el siguiente comando:

%%
%%\begin{figure}[h]
%%    \centering
%%    \includegraphics[width=0.85\textwidth]{figures/arquitectura.png}
%%    \caption{Arquitectura técnica preliminar del agente propuesto.}
%%    \label{fig:arquitectura}
%%\end{figure}
%%

\subsection{Soluciones de inteligencia artificial}

Explique qué técnicas de IA se utilizarán.

Ejemplos:

\begin{itemize}
    \item Modelos de lenguaje grandes.
    \item Sistemas basados en recuperación de información.
    \item Agentes conversacionales.
    \item Modelos de recomendación.
    \item Clasificadores o modelos predictivos.
\end{itemize}

Justifique la elección con base en literatura académica.

\subsection{Motor gráfico o plataforma de visualización}

Explique si usará una plataforma como Unity, Unreal Engine, WebGL, una aplicación web,
una aplicación móvil u otra tecnología.

\subsection{Servicios de voz}

Indique si el agente necesita:

\begin{itemize}
    \item Reconocimiento automático de voz.
    \item Síntesis de voz.
    \item Entrada por texto solamente.
    \item Interacción multimodal.
\end{itemize}

\subsection{Representación visual del agente}

Describa cómo se representa visualmente el agente.

Ejemplos:

\begin{itemize}
    \item Avatar 2D.
    \item Avatar 3D.
    \item Personaje humanoide.
    \item Mascota virtual.
    \item Interfaz sin cuerpo físico.
\end{itemize}

Explique por qué esa representación es adecuada para el problema y los usuarios.

\subsection{Entorno de despliegue}

Indique dónde funcionará el sistema:

\begin{itemize}
    \item PC.
    \item Móvil.
    \item Web.
    \item Realidad virtual.
    \item Realidad aumentada.
    \item Realidad extendida.
\end{itemize}

Justifique la elección según el contexto de uso.

\section{Conclusiones Preliminares}

Resuma brevemente:

\begin{itemize}
    \item El problema abordado.
    \item El agente propuesto.
    \item La relevancia de la investigación.
    \item Las principales brechas que se intentan abordar.
\end{itemize}

\printbibliography
\end{document}
